\documentclass[11pt]{article}

\usepackage[utf8]{inputenc}
\usepackage[T1]{fontenc}
\usepackage{amsmath,amssymb,amsfonts,bm}
\usepackage{geometry}
\geometry{margin=1in}

\title{Supplementary Material:\\
MacroMRT --- Instantaneous-Current Bayesian Inference\\
with State-Dependent Open-Channel Noise}
\author{}
\date{}

\begin{document}
\maketitle

\newcommand{\E}{\mathbb{E}}
\newcommand{\Var}{\operatorname{Var}}
\newcommand{\diag}{\operatorname{diag}}
\newcommand{\Nch}{N_{\mathrm{ch}}}
\newcommand{\gbar}{\overline{\gamma}_0}
\newcommand{\sbar}{\overline{\sigma^2}_0}

% ============================================================
\section{Setup and Notation}
% ============================================================

We consider $\Nch$ independent Markov ion channels with $K$ microscopic
states. At the start of an interval $[0,t]$ the macroscopic occupancy
distribution is summarised by a prior mean $\mu_0\in\mathbb{R}^{1\times K}$
(row vector of state probabilities) and a per-channel covariance
$\Sigma_0\in\mathbb{R}^{K\times K}$. Microscopic dynamics are governed
by the rate matrix $Q$ with transition matrix $P(t)=e^{Qt}$.

The recording produces an instantaneous-current measurement
$\overline{y}^{\mathrm{obs}}$ at the end of the interval. Per
boundary pair $(i_0,i_t)$ the model defines:
\begin{itemize}
  \item $\overline{\Gamma}_{i_0\to i_t}$: mean conductance over the
        recording window conditional on $X(0)=i_0,X(t)=i_t$;
  \item $\overline{V}_{i_0\to i_t}$: open-channel noise variance
        contribution conditional on $(i_0,i_t)$.
\end{itemize}

These boundary-conditioned moments are the same objects that MacroIR
uses; MacroMRT only requires their per-$i_0$ collapses
\begin{align}
  (\gbar)_{i_0}
  &= \sum_{i_t} P_{i_0\to i_t}(t)\,\overline{\Gamma}_{i_0\to i_t},
  \label{eq:gbar}\\
  (\sbar)_{i_0}
  &= \sum_{i_t} P_{i_0\to i_t}(t)\,\overline{V}_{i_0\to i_t}.
  \label{eq:sbar}
\end{align}
These are computed once per interval by the moment routine
(\texttt{Qdtm}). They reduce to the per-state constants
$\gamma_i,\sigma^2_i$ when $Q$ is constant on $[0,t]$ and the recording
window is short relative to kinetic timescales, but in general
\eqref{eq:gbar}--\eqref{eq:sbar} are required even for the M-family
because (i) within-window transitions modify the effective conductance
even at the endpoint, and (ii) agonist concentration changes within an
interval generate two distinct $Q$ matrices that must be folded through
the boundary kernel.

The instrument noise variance over the recording window is
$\epsilon^2_{0\to t}=\epsilon^2/t+\nu^2$ as in MacroIR.

\paragraph{Coordinates and singular covariance.}
The macroscopic occupancy $p$ lives on the simplex
$\{p\in\mathbb{R}^{1\times K}:p_i\ge 0,\sum_i p_i=1\}$; $\Sigma^{\mathrm{prop}}$ is
singular as a $K\times K$ matrix in ambient coordinates. We use
$(\Sigma^{\mathrm{prop}})^{-1}$ only as a formal device inside the
Sherman--Morrison/Woodbury algebra; all implementable formulas contain
only $\Sigma^{\mathrm{prop}}$.

% ============================================================
\section{Predictive Mean and Variance}
% ============================================================

The predictive mean is
\begin{equation}
  \overline{y}^{\mathrm{pred}}
  = \Nch\,\mu_0\cdot\gbar.
  \label{eq:ypred}
\end{equation}

The predictive variance is the sum of instrument noise and the
ensemble-averaged open-channel contribution:
\begin{equation}
  V
  = \epsilon^2_{0\to t}
  + \Nch\sum_{i_0}\mu_{0,i_0}\,(\sbar)_{i_0}.
  \label{eq:V}
\end{equation}

\paragraph{Comparison with MacroIR.}
MacroIR adds a third term, $\Nch\widetilde{\gamma^{T}\Sigma\gamma}$,
that captures the boundary cross-correlation between channel state and
current. MacroMRT \emph{drops} this term --- it is the central structural
difference between the M and I families. The MRT-vs-IRT comparison is
designed precisely to attribute the IR gain (if any) to this term and
the boundary-lifted cross-covariance, holding the σ² Taylor correction
fixed.

% ============================================================
\section{Energy, Gradient, Hessian}
% ============================================================

Let $p$ denote a candidate posterior occupancy. Define
\begin{align}
  \delta(p) &= \overline{y}^{\mathrm{obs}} - \Nch(p\cdot\gbar),\\
  V(p)      &= \epsilon^2_{0\to t} + \Nch(p\cdot\sbar).
\end{align}
The Gaussian negative log-likelihood is
\[
  F_{\mathrm{like}}(p)
  = \tfrac12\log V(p) + \tfrac12\delta(p)^2/V(p)
  \quad(+\,\text{const}).
\]
The Gaussian negative log-prior over $p$, with precision
$\Lambda_p=\Nch(\Sigma^{\mathrm{prop}})^{-1}$, is
\[
  F_{\mathrm{prior}}(p)
  = \tfrac12(p-\mu^{\mathrm{prior}})\Lambda_p(p-\mu^{\mathrm{prior}})^{T}.
\]
Define the energy $E(p):=2(F_{\mathrm{like}}+F_{\mathrm{prior}})$ to
absorb the $\tfrac12$ factors. Up to additive constants
\begin{equation}
  E(p) = \log V(p) + \frac{\delta(p)^2}{V(p)}
  + (p-\mu^{\mathrm{prior}})\Nch(\Sigma^{\mathrm{prop}})^{-1}(p-\mu^{\mathrm{prior}})^{T}.
  \label{eq:E}
\end{equation}

The MAP and Laplace covariance satisfy
$\Sigma_{\mathrm{Lap}}=2(\nabla^2 E)^{-1}$ at the MAP.

\subsection{Tilted vector and gradient}

Following MacroTaylor, define
\begin{equation}
  v(p) = \gbar + \frac{\delta(p)}{V(p)}\,\sbar.
  \label{eq:vdef}
\end{equation}
With $\nabla\delta=-\Nch\gbar^{T}$, $\nabla V=\Nch\sbar^{T}$,
\begin{equation}
  \mathbf{g}(p) = \nabla E(p)
  = 2\Nch(p-\mu^{\mathrm{prior}})(\Sigma^{\mathrm{prop}})^{-1}
  - \Nch\!\left[\frac{2\delta}{V}\gbar^{T}
                  +\Big(\frac{\delta^2}{V^2}-\frac{1}{V}\Big)\sbar^{T}\right].
  \label{eq:grad}
\end{equation}
At $p=\mu^{\mathrm{prior}}$ the prior term vanishes and
$\mathbf{g}=-\Nch\frac{\delta}{V}(\gbar^{T}+v^{T})$.

\subsection{Hessian: rank-2 structure}

Direct differentiation gives, with $\nabla^2\delta=\nabla^2 V=0$,
\begin{align}
  \nabla^2(\log V)
  &= -\frac{\Nch^2}{V^2}\,\sbar\sbar^{T},\\
  \nabla^2(\delta^2/V)
  &= 2\Nch^2\!\left[\frac{1}{V}\gbar\gbar^{T}
        + \frac{\delta}{V^2}(\gbar\sbar^{T}+\sbar\gbar^{T})
        + \frac{\delta^2}{V^3}\sbar\sbar^{T}\right].
\end{align}
Adding the prior Hessian $2\Nch(\Sigma^{\mathrm{prop}})^{-1}$ and
combining:
\begin{equation}
  \mathbf{H}(p)
  = 2\Nch\,(\Sigma^{\mathrm{prop}})^{-1}
  + 2\Nch^2\!\left[\frac{1}{V}\gbar\gbar^{T}
       + \frac{\delta}{V^2}(\gbar\sbar^{T}+\sbar\gbar^{T})
       + \Big(\frac{\delta^2}{V^3}-\frac{1}{2V^2}\Big)\sbar\sbar^{T}\right].
  \label{eq:hess}
\end{equation}
Setting $u_1=v=\gbar+\frac{\delta}{V}\sbar$, $u_2=\sbar$ and using the
identity $\frac{\Nch}{V}vv^{T}$ to absorb the cross terms,
\eqref{eq:hess} becomes the rank-2 update of the prior precision
\begin{equation}
  \mathbf{H}(p)
  = 2\Nch\bigl[(\Sigma^{\mathrm{prop}})^{-1} + UCU^{T}\bigr],
  \quad
  U=[v,\sbar],
  \quad
  C=\diag\!\left(\frac{\Nch}{V},\,-\frac{\Nch}{2V^2}\right).
  \label{eq:rank2}
\end{equation}

% ============================================================
\section{Rank-1 Quasi-Laplace Scheme (Main Implementation)}
% ============================================================

Dropping the $\log V$ likelihood term gives the quasi-likelihood energy
$E_{\mathrm{quasi}}=\delta^2/V+P$ whose Hessian is the rank-1 update
\begin{equation}
  \mathbf{H}_{\mathrm{quasi}}(p)
  = 2\Nch\bigl[(\Sigma^{\mathrm{prop}})^{-1} + (\Nch/V)\,vv^{T}\bigr].
  \label{eq:rank1}
\end{equation}
Sherman--Morrison gives, with $s=v^{T}\Sigma^{\mathrm{prop}}v$,
\begin{equation}
  \boxed{
    \Sigma^{\mathrm{post}}
    = \Sigma^{\mathrm{prop}}
    - \frac{\Nch}{V+\Nch s}\,
      \Sigma^{\mathrm{prop}}vv^{T}\Sigma^{\mathrm{prop}}.
  }
  \label{eq:Sigma_post}
\end{equation}
The mean update at $p_0=\mu^{\mathrm{prior}}$ is
\begin{equation}
  \boxed{
    \mu^{\mathrm{post}}
    = \mu^{\mathrm{prior}}
    + \frac{\delta}{2V}\,\Sigma^{\mathrm{post}}\,(\gbar+v),
  }
  \label{eq:mu_post}
\end{equation}
which is the Newton step from the prior mean using the rank-1
Hessian. No $\Sigma_p^{-1}$ survives in the implementable expressions
\eqref{eq:Sigma_post}--\eqref{eq:mu_post}.

\paragraph{Working scheme.}
For ensembles of $\Nch\sim 10^2$--$10^4$ channels the quasi-likelihood
update is the operational form: the $\log V$ correction is $O(1/\Nch)$
relative to the leading data term and is dominated by other modelling
errors. The rank-2 form is retained for completeness in
Section~\ref{sec:rank2-exact}.

% ============================================================
\section{Rank-2 Exact Laplace Scheme}\label{sec:rank2-exact}
% ============================================================

For small $\Nch$ or stress tests, the rank-2 update from
\eqref{eq:rank2} requires the Woodbury identity. Set
\begin{equation}
  K = \bigl(C^{-1} + U^{T}\Sigma^{\mathrm{prop}}U\bigr)^{-1}
  \in\mathbb{R}^{2\times 2},
  \label{eq:Kdef}
\end{equation}
i.e.\ explicitly
\[
  K^{-1} =
  \begin{pmatrix}
    V/\Nch + s & b\\
    b          & -2V^2/\Nch + c
  \end{pmatrix},
  \quad
  s=v^{T}\Sigma^{\mathrm{prop}}v,
  \quad
  b=v^{T}\Sigma^{\mathrm{prop}}\sbar,
  \quad
  c=\sbar^{T}\Sigma^{\mathrm{prop}}\sbar.
\]
The Laplace covariance is
\begin{equation}
  \Sigma^{\mathrm{post}}_{\mathrm{exact}}(p_0)
  = \Sigma^{\mathrm{prop}}
  - \Sigma^{\mathrm{prop}}U(p_0)K(p_0)U(p_0)^{T}\Sigma^{\mathrm{prop}}.
  \label{eq:Sigma_post_exact}
\end{equation}
Eliminating the formal $(\Sigma^{\mathrm{prop}})^{-1}$ in the Newton
step (cf.\ MacroTaylor supplement, §3.2) yields the mean update
\begin{equation}
  \mu^{\mathrm{post}}_{\mathrm{exact}}
  = \mu^{\mathrm{prior}}
  + \Delta p\,U K U^{T}\Sigma^{\mathrm{prop}}
  + \frac{\Nch}{2}\,q^{T}\Sigma^{\mathrm{post}}_{\mathrm{exact}}(p_0),
  \label{eq:mu_post_exact}
\end{equation}
with $\Delta p=p_0-\mu^{\mathrm{prior}}$ and
$q^{T}=(2\delta/V)\gbar^{T}+(\delta^2/V^2-1/V)\sbar^{T}$.

% ============================================================
\section{Newton-MAP Iteration}
% ============================================================

For the rank-1 scheme the standard recipe is:
\begin{enumerate}
  \item Initialise $p^{(0)}=\mu^{\mathrm{prior}}$.
  \item For $k=0,1,\dots$ until convergence:
  \begin{enumerate}
    \item Compute $\delta^{(k)},V^{(k)},v^{(k)},s^{(k)}$ from $p^{(k)}$.
    \item Form $\Sigma^{(k)}_{\mathrm{post}}$ via \eqref{eq:Sigma_post}.
    \item Update via the general step
    \[
      p^{(k+1)} = \mu^{\mathrm{prior}}
      + \frac{\Nch}{V^{(k)}+\Nch s^{(k)}}\,\Delta p^{(k)} v^{(k)}v^{(k)\,T}\Sigma^{\mathrm{prop}}
      + \frac{\Nch\delta^{(k)}}{V^{(k)}}\,\gbar^{T}\Sigma^{(k)}_{\mathrm{post}}
      + \frac{\Nch\delta^{(k)\,2}}{2V^{(k)\,2}}\,\sbar^{T}\Sigma^{(k)}_{\mathrm{post}}.
    \]
    \item Project $p^{(k+1)}$ to the simplex (probability floor, cf.\
          MacroIR variance-inflation correction).
  \end{enumerate}
\end{enumerate}
A single iteration from $p^{(0)}=\mu^{\mathrm{prior}}$ recovers
\eqref{eq:mu_post}; in the M-family one iteration is typically used
per measurement.

% ============================================================
\section{Reduction Lemmas}
% ============================================================

\subsection{$\sbar=0$ recovers MacroMR (Moffatt 2007)}

Setting $\sbar\equiv 0$ gives $v=\gbar$, $s=\gbar^{T}\Sigma^{\mathrm{prop}}\gbar$,
$V=\epsilon^2_{0\to t}$, and
\begin{align}
  \Sigma^{\mathrm{post}}_{\mathrm{MR}}
  &= \Sigma^{\mathrm{prop}} - \frac{\Nch}{V+\Nch\,\gbar^{T}\Sigma^{\mathrm{prop}}\gbar}
     \Sigma^{\mathrm{prop}}\gbar\gbar^{T}\Sigma^{\mathrm{prop}},\\
  \mu^{\mathrm{post}}_{\mathrm{MR}}
  &= \mu^{\mathrm{prior}} + \frac{\delta}{V}\,\Sigma^{\mathrm{post}}_{\mathrm{MR}}\gbar.
\end{align}
This is the Moffatt 2007 macroscopic recursive update written in the
boundary-conditioned conventions used throughout MacroIR/MacroIRT.

\subsection{Tilde-collapse to MacroIR / MacroIRT}

The IR family lifts the cross-covariance $\Sigma^{\mathrm{prop}}\gbar$
to the boundary tilde $\widetilde{\gamma^{T}\Sigma}$, the scalar
$v^{T}\Sigma^{\mathrm{prop}}v$ to $\widetilde{v^{T}\Sigma v}$, and adds
the third predictive variance term $\Nch\widetilde{\gamma^{T}\Sigma\gamma}$.
Substituting these replacements in
\eqref{eq:Sigma_post}--\eqref{eq:mu_post} gives the MacroIRT update
exactly. MacroIR is recovered by additionally setting $\sbar=0$,
$v=\gbar$, $s=\widetilde{\gamma^{T}\Sigma\gamma}$.

% ============================================================
\section{Trust-Region Simplex Shrink}
% ============================================================

The Kalman step \eqref{eq:mu_post} can push components of
$\mu^{\mathrm{post}}$ outside $[p_{\min},p_{\max}]$. As in MacroIR
(cf.\ \texttt{macro\_ir\_variance\_inflation\_correction.tex}), we
introduce a shrink factor $\alpha^\star\in[0,1]$ chosen as the largest
scalar that keeps every component inside the simplex floor, applied
uniformly to both the mean and covariance updates. The likelihood
$\log L$ is computed with the unmodified $V$; $\alpha^\star$ is
recorded as a diagnostic (\texttt{trust\_coefficient}). The mechanism
and rationale are identical to MacroIR's; see that document for the
full derivation and discussion of the discarded variance-inflation
interpretation.

% ============================================================
\section{Log-Likelihood Contribution}
% ============================================================

\begin{equation}
  \log L = -\tfrac{1}{2}\bigl[\log(2\pi V) + \delta^2/V\bigr],
  \label{eq:logL}
\end{equation}
identical in form to MacroIR's per-interval scalar Gaussian
log-likelihood, with the MRT-specific $V$ from \eqref{eq:V}. The
parameter expectation $\mathrm{E}[\log L]$ over the prior follows the
same scalar-Gaussian derivation as MacroIR.

\end{document}
